\clearpage\pagestyle{otMain}\fontsize{11}{15}\selectfont\corrigesChapitre{6}{Proportionnalité, ratios et pourcentages}
\begin{blocLivre}{cor-6-1}\begin{corrige}{1}{Passer par l’unité}\demarche{Calculer le prix d’une unité avant celui du lot.}
Un cahier coûte $10{,}80/4=2{,}70$ €. Sept coûtent $7\times2{,}70=\res{18{,}90\text{ €}}$.
\end{corrige}\end{blocLivre}
\begin{blocLivre}{cor-6-2}\begin{corrige}{2}{Un tableau proportionnel ?}\demarche{Comparer tous les rapports correspondant à la même direction.}
Les trois rapports prix sur quantité valent $7/2=17{,}5/5=28/8=3{,}5$. Le tableau est proportionnel, de coefficient $\res{3{,}5}$, et modélisé par $\res{x\mapsto3{,}5x}$.
\end{corrige}\end{blocLivre}
\begin{blocLivre}{cor-6-3}\begin{corrige}{3}{Appliquer une réduction}\demarche{Multiplier par le coefficient de baisse.}
Le coefficient est $1-25/100=0{,}75$. Le nouveau prix vaut $64\times0{,}75=\res{48\text{ €}}$. La réduction est de 16 €.
\end{corrige}\end{blocLivre}
\begin{blocLivre}{cor-6-4}\begin{corrige}{4}{Calculer un taux}\demarche{Diviser la variation par la valeur de départ.}
La hausse vaut $184-160=24$. Rapportée à l’effectif initial, elle représente $24/160=0{,}15$, soit $\res{15\,\%}$.
\end{corrige}\end{blocLivre}
\begin{blocLivre}{cor-6-5}\begin{corrige}{5}{Partager selon un ratio}\demarche{Compter toutes les parts du ratio.}
Il y a $3+5=8$ parts, chacune de $56/8=7$ billes. Les deux groupes reçoivent $\res{21\text{ et }35\text{ billes}}$. Le total et le ratio sont vérifiés.
\end{corrige}\end{blocLivre}
\begin{blocLivre}{cor-6-6}\begin{corrige}{6}{Retrouver un prix initial}\demarche{Diviser par le coefficient d’évolution.}
Si le prix initial est $P$, $1{,}08P=81$. Ainsi $P=81/1{,}08=\res{75\text{ €}}$. Soustraire 8 \% de 81 utiliserait une autre référence.
\end{corrige}\end{blocLivre}
\begin{blocLivre}{cor-6-7}\begin{corrige}{7}{Deux baisses}\demarche{Appliquer chaque taux à la valeur actualisée.}
Les coefficients se multiplient : $250\times0{,}9\times0{,}8=\res{180\text{ €}}$. Le coefficient global est 0,72, donc la baisse totale est $\res{28\,\%}$, et non 30 \%.
\end{corrige}\end{blocLivre}
\begin{blocLivre}{cor-6-8}\begin{corrige}{8}{Une échelle}\demarche{Mettre les longueurs dans la même unité.}
On convertit 7,2 km en 720000 cm. Le rapport plan sur réel vaut $9/720000=1/80000$ : l’échelle est $\res{1/80000}$. Ainsi 4 cm représentent $4\times80000=320000$ cm, soit $\res{3{,}2\text{ km}}$.
\end{corrige}\end{blocLivre}
\begin{blocLivre}{cor-6-9}\begin{corrige}{9}{Une masse en trois parts}\demarche{Additionner les termes du ratio avant de partager.}
Les neuf parts totalisent 4,5 kg ; une part vaut 0,5 kg. Les masses sont donc $\res{2{,}5\text{ kg},\ 1{,}5\text{ kg},\ 0{,}5\text{ kg}}$. Leur somme vaut bien 4,5 kg.
\end{corrige}\end{blocLivre}
\begin{blocLivre}{cor-6-10}\begin{corrige}{10}{Proportionnellement à l’âge}\demarche{Simplifier le ratio, puis partager le total.}
Le partage suit le ratio $12:18$, soit $2:3$. Il y a 5 parts de $450/5=90$ €. Les sœurs reçoivent $\res{180\text{ €}}$ et $\res{270\text{ €}}$. Contrôle : $180/12=270/18=15$ € par année d’âge.
\end{corrige}\end{blocLivre}
\begin{blocLivre}{cor-6-11}\begin{corrige}{11}{Comparer deux promotions}\demarche{Traduire exactement l’ordre des avantages.}
A conduit à $120-30=\res{90\text{ €}}$. B conduit à $120\times0{,}8-10=\res{86\text{ €}}$. Avec les conditions indiquées, B fait payer 4 € de moins. Le calcul dépend de l’utilisation immédiate du bon.
\end{corrige}\end{blocLivre}
\begin{blocLivre}{cor-6-12}\begin{corrige}{12}{Un mélange à corriger}\demarche{Fixer la quantité conservée et chercher le nouveau total.}
Les 2 L de sirop doivent représenter 0,16 du volume final $V$ : $V=2/0{,}16=12{,}5$ L. Le mélange initial fait 10 L. Il faut ajouter $\res{2{,}5\text{ L d’eau}}$. Le volume de sirop reste 2 L, mais le total augmente.
\end{corrige}\end{blocLivre}
\begin{blocLivre}{cor-6-13}\begin{corrige}{13}{Deux pourcentages emboîtés}\demarche{Nommer le groupe de référence de chaque proportion.}
L’association compte $720\times0{,}45=324$ élèves. Le badminton en compte $324\times2/9=\res{72}$. Leur part dans le collège est $72/720=\res{10\,\%}$. On multiplie 45 \% par $2/9$, pas par 2.
\end{corrige}\end{blocLivre}
\begin{blocLivre}{cor-6-14}\begin{corrige}{14}{Lire une annonce}\demarche{Distinguer points de pourcentage, taux relatif et effectifs.}
La part augmente de $\res{3\text{ points}}$. Relativement à 12 \%, la hausse vaut $(15-12)/12=\res{25\,\%}$. Les effectifs passent de 96 à 120, soit $\res{24}$ cyclistes supplémentaires. Le maintien du total de 800 est indispensable à cette dernière comparaison.
\end{corrige}\end{blocLivre}
\begin{blocLivre}{cor-6-15}\begin{corrige}{15}{Proportionnalité et alignement}\demarche{Repérer deux triangles emboîtés de sommet $O$.}
Les droites $(AH)$ et $(BK)$ sont toutes deux perpendiculaires à l’axe des abscisses : elles sont parallèles. Les points $O,H,K$ et $O,A,B$ sont alignés. Thalès donne $\dfrac{AH}{BK}=\dfrac{OH}{OK}$, soit $\dfrac{3}{7{,}5}=\dfrac{4}{10}$. Donc $\dfrac34=\dfrac{7{,}5}{10}=\res{0{,}75}$ : les ordonnées s’obtiennent en multipliant les abscisses par 0,75.
\label{fig-cor-6-15}\figureGeo[corrige]{thales-repere}
\end{corrige}\end{blocLivre}
